\subsection{Two universal characteristic comparisons}

\begin{statement}{P3R.1}{Proved}
Let $W$ be a compact almost complex eight-manifold whose
outward-oriented boundary is the homotopy sphere $\Sigma$.
In the presentation $\Theta_7=\Z/28$ fixed on the first page,
\[
 [\Sigma]=
 \frac{-c_1^4+4c_1^2c_2+3c_2^2-3\sigma(W)}{24}
 \pmod {28}.
\]
All characteristic products are evaluated on $[W,\partial W]$
using their unique relative rational lifts. A spin assumption
on $W$ is unnecessary.
\end{statement}
\begin{proof}
The almost complex structure induces a spin-c structure
with determinant class $c_1=c_1(TW)$ and the complex
orientation. The homotopy sphere has $H^1=H^2=H^3=H^4=0$
over $\Q$ and a unique spin structure. Thus the boundary
determinant is trivial, the induced boundary spin structure
is the required one, and the exact sequences of the pair
give unique relative rational lifts in degrees two and four.
The spin-c Eells--Kuiper formula in \cite[Section~3.1(v)]{Jiang2014},
with these hypotheses, is
\[
 \mu(\Sigma)=
 \frac{p_1^2-4\sigma(W)}{896}
 -\frac{c_1^2p_1}{192}+\frac{c_1^4}{384}
 \pmod\Z.
\]
The formula is for the outward boundary, as used here.
For the underlying real tangent bundle of a complex bundle,
$p_1=c_1^2-2c_2$. Multiplication by $28$ and substitution
give
\[
 \begin{split}
 28\mu(\Sigma)
 &=\frac{3p_1^2-14c_1^2p_1+7c_1^4-12\sigma(W)}{96}\\
 &=\frac{-c_1^4+4c_1^2c_2+3c_2^2-3\sigma(W)}{24}
 \pmod {28}.
 \end{split}
\]
This is precisely the specified presentation of $\Theta_7$.
The substitution is independently executed in
\texttt{code/root\_comparisons/characteristics.py}.
\end{proof}

\begin{statement}{P3R.2}{Proved}
The equation $6a=b$ in $\Z/28$ has no solution for odd $b$.
For $b=2r$, its solutions are exactly
$a=5r$ and $a=5r+14$ modulo $28$.
Thus the sixfold class alone leaves precisely these two
possibilities for the original class; absent a value of $b$,
this calculation excludes no original class.
\end{statement}
\begin{proof}
Reduction modulo two gives the first assertion. For even
$b$, division by two identifies the equation with
$3a=r$ in $\Z/14$. The inverse of $3$ in that ring is
$5$, since $3\cdot5=1+14$. Its unique solution modulo
$14$ has the two displayed lifts modulo $28$.
Conversely each lift satisfies the original equation.
Every original class occurs for its own value $b=6a$.
The attached program verifies the complete table of
twenty-eight target residues. The integral index used
in P3K.7--8 supplies the additional information that
settles the actual original family.
\end{proof}
